\documentclass[aspectratio=169,10pt]{beamer}

% ═════════════════════════════════════════════════════════════════════════════
%  MINIMAL ACADEMIC BEAMER THEME — CvG-Diff
%  Philosophy: Content flows naturally; boxes only for key takeaways.
%  Based on:   Metropolis (structure) + custom lightweight styling
% ═════════════════════════════════════════════════════════════════════════════

% ─── Core Theme ─────────────────────────────────────────────────────────────
\usetheme[progressbar=frametitle,block=fill,numbering=fraction]{metropolis}
\usepackage{appendixnumberbeamer}

% ─── Fonts ──────────────────────────────────────────────────────────────────
\usepackage{unicode-math}
\setmainfont{Latin Modern Roman}
\setsansfont{Latin Modern Sans}
\setmathfont{Latin Modern Math}
\usepackage{amsmath,amssymb}

% ─── CJK Support (cross-platform) ───────────────────────────────────────────
% ============================================================================
% font-config.sty — Cross-Platform CJK Font Configuration
% ============================================================================
% Single source of truth for CJK font setup across all platforms.
%
% USER OVERRIDE (highest priority):
%   Before loading this package, define:
%     \def\CJKFontPath{/path/to/fonts/}
%     \def\CJKFontName{FontFile.ttf}
%     \def\CJKFontBold{FontFileBold.ttf}   (optional)
%
%   Example for agent chat:
%     \def\CJKFontPath{/usr/share/fonts/truetype/noto/}
%     \def\CJKFontName{NotoSansSC-Regular.ttf}
%     \def\CJKFontBold{NotoSansSC-Bold.ttf}
%     % ============================================================================
% font-config.sty — Cross-Platform CJK Font Configuration
% ============================================================================
% Single source of truth for CJK font setup across all platforms.
%
% USER OVERRIDE (highest priority):
%   Before loading this package, define:
%     \def\CJKFontPath{/path/to/fonts/}
%     \def\CJKFontName{FontFile.ttf}
%     \def\CJKFontBold{FontFileBold.ttf}   (optional)
%
%   Example for agent chat:
%     \def\CJKFontPath{/usr/share/fonts/truetype/noto/}
%     \def\CJKFontName{NotoSansSC-Regular.ttf}
%     \def\CJKFontBold{NotoSansSC-Bold.ttf}
%     % ============================================================================
% font-config.sty — Cross-Platform CJK Font Configuration
% ============================================================================
% Single source of truth for CJK font setup across all platforms.
%
% USER OVERRIDE (highest priority):
%   Before loading this package, define:
%     \def\CJKFontPath{/path/to/fonts/}
%     \def\CJKFontName{FontFile.ttf}
%     \def\CJKFontBold{FontFileBold.ttf}   (optional)
%
%   Example for agent chat:
%     \def\CJKFontPath{/usr/share/fonts/truetype/noto/}
%     \def\CJKFontName{NotoSansSC-Regular.ttf}
%     \def\CJKFontBold{NotoSansSC-Bold.ttf}
%     \input{font-config.sty}
%
% AUTO-DETECT priority:
%   1. User override (\CJKFontPath defined)
%   2. Windows: Microsoft YaHei (C:/WINDOWS/Fonts/)
%   3. Linux: Noto Sans CJK (/usr/share/fonts/...)
%   4. macOS: Noto Sans CJK (~/Library/Fonts/ or /System/Library/Fonts/)
%   5. Generic: Noto Sans SC in common paths
%
% REQUIREMENTS:
%   - XeLaTeX (fontspec + xeCJK)
%   - CJK fonts installed by user (we do NOT distribute fonts)
% ============================================================================

\NeedsTeXFormat{LaTeX2e}
\ProvidesPackage{font-config}[2026/05/23 Cross-Platform CJK Font Config]

\RequirePackage{fontspec}
\RequirePackage{xeCJK}

% ── Internal helpers ────────────────────────────────────────────────────────
% Use \makeatletter/\makeatother because this file may be \input'd
% (not loaded as a proper .sty package), where @ is not a letter.
\makeatletter
\newcommand{\@fc@setfonts}[4][]{%
  % #1 = optional RawFeature for variable fonts
  % #2 = font path (with trailing slash)
  % #3 = main font file
  % #4 = bold font file (or empty)
  \def\@fc@arg{#4}%
  \ifx\@fc@arg\empty
    % No separate bold font
    \setCJKmainfont{#3}[Path=#2,#1]%
    \setCJKsansfont{#3}[Path=#2,#1]%
    \setCJKmonofont{#3}[Path=#2,#1]%
  \else
    \setCJKmainfont{#3}[Path=#2,BoldFont=#4]%
    \setCJKsansfont{#3}[Path=#2,BoldFont=#4]%
    \setCJKmonofont{#3}[Path=#2,BoldFont=#4]%
  \fi
}

\newcommand{\@fc@warn}[1]{%
  \PackageWarningNoLine{font-config}{#1}%
}

% ── Priority 1: User Override ───────────────────────────────────────────────
\ifx\CJKFontPath\undefined
  % No user override — proceed to auto-detect
\else
  \@fc@warn{Using user-defined CJK font: \CJKFontPath\CJKFontName}
  \ifx\CJKFontBold\undefined
    \@fc@setfonts{\CJKFontPath}{\CJKFontName}{}%
  \else
    \@fc@setfonts{\CJKFontPath}{\CJKFontName}{\CJKFontBold}%
  \fi
  % Skip rest of file
  \expandafter\endinput
\fi

% ── Priority 1.5: Project-local .fonts/ — Source Han Serif SC ────────────────
% Ships as individual OTFs in .fonts/ (gitignored). Prefer Medium weight for
% body text on slides — heavier than the system .ttc Regular.
\IfFileExists{./.fonts/SourceHanSerifSC-Medium.otf}{%
  \@fc@warn{Detected project-local Source Han Serif SC (Medium+Bold)}
  \@fc@setfonts{./.fonts/}{SourceHanSerifSC-Medium.otf}{SourceHanSerifSC-Bold.otf}%
  \expandafter\endinput
}{}

% ── Priority 2: Windows — Microsoft YaHei ───────────────────────────────────
\IfFileExists{C:/WINDOWS/Fonts/msyh.ttc}{%
  \@fc@warn{Detected Windows — using Microsoft YaHei}
  \@fc@setfonts{C:/WINDOWS/Fonts/}{msyh.ttc}{msyhbd.ttc}%
}{%
% ── Priority 3: Linux — Noto Sans CJK (system package) ──────────────────────
  \IfFileExists{/usr/share/fonts/opentype/noto/NotoSansCJK-Regular.ttc}{%
    \@fc@warn{Detected Linux — using Noto Sans CJK (opentype)}
    \@fc@setfonts{/usr/share/fonts/opentype/noto/}{NotoSansCJK-Regular.ttc}{NotoSansCJK-Bold.ttc}%
  }{%
    \IfFileExists{/usr/share/fonts/truetype/noto/NotoSansSC-Regular.ttf}{%
      \@fc@warn{Detected Linux — using Noto Sans SC (truetype)}
      \@fc@setfonts{/usr/share/fonts/truetype/noto/}{NotoSansSC-Regular.ttf}{NotoSansSC-Bold.ttf}%
    }{%
% ── Priority 4: macOS — Noto Sans CJK or system fonts ───────────────────────
      \IfFileExists{/System/Library/Fonts/PingFang.ttc}{%
        \@fc@warn{Detected macOS — using PingFang}
        \@fc@setfonts{/System/Library/Fonts/}{PingFang.ttc}{}%
      }{%
        \IfFileExists{~/Library/Fonts/NotoSansSC-Regular.ttf}{%
          \@fc@warn{Detected macOS — using user-installed Noto Sans SC}
          \@fc@setfonts{~/Library/Fonts/}{NotoSansSC-Regular.ttf}{NotoSansSC-Bold.ttf}%
        }{%
% ── Priority 5: Generic fallback — search common paths ──────────────────────
          \IfFileExists{/usr/local/share/fonts/NotoSansSC-Regular.ttf}{%
            \@fc@warn{Using generic Noto Sans SC}
            \@fc@setfonts{/usr/local/share/fonts/}{NotoSansSC-Regular.ttf}{NotoSansSC-Bold.ttf}%
          }{%
% ── FAILURE: No CJK font found ──────────────────────────────────────────────
            \@fc@warn{No CJK font found on system. CJK text will not render correctly.}
            \@fc@warn{Install fonts: sudo apt-get install fonts-noto-cjk  (Linux)}
            \@fc@warn{Or set manually: \def\CJKFontPath{/path/} \def\CJKFontName{font.ttf}}
          }%
        }%
      }%
    }%
  }%
}
\makeatother

%
% AUTO-DETECT priority:
%   1. User override (\CJKFontPath defined)
%   2. Windows: Microsoft YaHei (C:/WINDOWS/Fonts/)
%   3. Linux: Noto Sans CJK (/usr/share/fonts/...)
%   4. macOS: Noto Sans CJK (~/Library/Fonts/ or /System/Library/Fonts/)
%   5. Generic: Noto Sans SC in common paths
%
% REQUIREMENTS:
%   - XeLaTeX (fontspec + xeCJK)
%   - CJK fonts installed by user (we do NOT distribute fonts)
% ============================================================================

\NeedsTeXFormat{LaTeX2e}
\ProvidesPackage{font-config}[2026/05/23 Cross-Platform CJK Font Config]

\RequirePackage{fontspec}
\RequirePackage{xeCJK}

% ── Internal helpers ────────────────────────────────────────────────────────
% Use \makeatletter/\makeatother because this file may be \input'd
% (not loaded as a proper .sty package), where @ is not a letter.
\makeatletter
\newcommand{\@fc@setfonts}[4][]{%
  % #1 = optional RawFeature for variable fonts
  % #2 = font path (with trailing slash)
  % #3 = main font file
  % #4 = bold font file (or empty)
  \def\@fc@arg{#4}%
  \ifx\@fc@arg\empty
    % No separate bold font
    \setCJKmainfont{#3}[Path=#2,#1]%
    \setCJKsansfont{#3}[Path=#2,#1]%
    \setCJKmonofont{#3}[Path=#2,#1]%
  \else
    \setCJKmainfont{#3}[Path=#2,BoldFont=#4]%
    \setCJKsansfont{#3}[Path=#2,BoldFont=#4]%
    \setCJKmonofont{#3}[Path=#2,BoldFont=#4]%
  \fi
}

\newcommand{\@fc@warn}[1]{%
  \PackageWarningNoLine{font-config}{#1}%
}

% ── Priority 1: User Override ───────────────────────────────────────────────
\ifx\CJKFontPath\undefined
  % No user override — proceed to auto-detect
\else
  \@fc@warn{Using user-defined CJK font: \CJKFontPath\CJKFontName}
  \ifx\CJKFontBold\undefined
    \@fc@setfonts{\CJKFontPath}{\CJKFontName}{}%
  \else
    \@fc@setfonts{\CJKFontPath}{\CJKFontName}{\CJKFontBold}%
  \fi
  % Skip rest of file
  \expandafter\endinput
\fi

% ── Priority 1.5: Project-local .fonts/ — Source Han Serif SC ────────────────
% Ships as individual OTFs in .fonts/ (gitignored). Prefer Medium weight for
% body text on slides — heavier than the system .ttc Regular.
\IfFileExists{./.fonts/SourceHanSerifSC-Medium.otf}{%
  \@fc@warn{Detected project-local Source Han Serif SC (Medium+Bold)}
  \@fc@setfonts{./.fonts/}{SourceHanSerifSC-Medium.otf}{SourceHanSerifSC-Bold.otf}%
  \expandafter\endinput
}{}

% ── Priority 2: Windows — Microsoft YaHei ───────────────────────────────────
\IfFileExists{C:/WINDOWS/Fonts/msyh.ttc}{%
  \@fc@warn{Detected Windows — using Microsoft YaHei}
  \@fc@setfonts{C:/WINDOWS/Fonts/}{msyh.ttc}{msyhbd.ttc}%
}{%
% ── Priority 3: Linux — Noto Sans CJK (system package) ──────────────────────
  \IfFileExists{/usr/share/fonts/opentype/noto/NotoSansCJK-Regular.ttc}{%
    \@fc@warn{Detected Linux — using Noto Sans CJK (opentype)}
    \@fc@setfonts{/usr/share/fonts/opentype/noto/}{NotoSansCJK-Regular.ttc}{NotoSansCJK-Bold.ttc}%
  }{%
    \IfFileExists{/usr/share/fonts/truetype/noto/NotoSansSC-Regular.ttf}{%
      \@fc@warn{Detected Linux — using Noto Sans SC (truetype)}
      \@fc@setfonts{/usr/share/fonts/truetype/noto/}{NotoSansSC-Regular.ttf}{NotoSansSC-Bold.ttf}%
    }{%
% ── Priority 4: macOS — Noto Sans CJK or system fonts ───────────────────────
      \IfFileExists{/System/Library/Fonts/PingFang.ttc}{%
        \@fc@warn{Detected macOS — using PingFang}
        \@fc@setfonts{/System/Library/Fonts/}{PingFang.ttc}{}%
      }{%
        \IfFileExists{~/Library/Fonts/NotoSansSC-Regular.ttf}{%
          \@fc@warn{Detected macOS — using user-installed Noto Sans SC}
          \@fc@setfonts{~/Library/Fonts/}{NotoSansSC-Regular.ttf}{NotoSansSC-Bold.ttf}%
        }{%
% ── Priority 5: Generic fallback — search common paths ──────────────────────
          \IfFileExists{/usr/local/share/fonts/NotoSansSC-Regular.ttf}{%
            \@fc@warn{Using generic Noto Sans SC}
            \@fc@setfonts{/usr/local/share/fonts/}{NotoSansSC-Regular.ttf}{NotoSansSC-Bold.ttf}%
          }{%
% ── FAILURE: No CJK font found ──────────────────────────────────────────────
            \@fc@warn{No CJK font found on system. CJK text will not render correctly.}
            \@fc@warn{Install fonts: sudo apt-get install fonts-noto-cjk  (Linux)}
            \@fc@warn{Or set manually: \def\CJKFontPath{/path/} \def\CJKFontName{font.ttf}}
          }%
        }%
      }%
    }%
  }%
}
\makeatother

%
% AUTO-DETECT priority:
%   1. User override (\CJKFontPath defined)
%   2. Windows: Microsoft YaHei (C:/WINDOWS/Fonts/)
%   3. Linux: Noto Sans CJK (/usr/share/fonts/...)
%   4. macOS: Noto Sans CJK (~/Library/Fonts/ or /System/Library/Fonts/)
%   5. Generic: Noto Sans SC in common paths
%
% REQUIREMENTS:
%   - XeLaTeX (fontspec + xeCJK)
%   - CJK fonts installed by user (we do NOT distribute fonts)
% ============================================================================

\NeedsTeXFormat{LaTeX2e}
\ProvidesPackage{font-config}[2026/05/23 Cross-Platform CJK Font Config]

\RequirePackage{fontspec}
\RequirePackage{xeCJK}

% ── Internal helpers ────────────────────────────────────────────────────────
% Use \makeatletter/\makeatother because this file may be \input'd
% (not loaded as a proper .sty package), where @ is not a letter.
\makeatletter
\newcommand{\@fc@setfonts}[4][]{%
  % #1 = optional RawFeature for variable fonts
  % #2 = font path (with trailing slash)
  % #3 = main font file
  % #4 = bold font file (or empty)
  \def\@fc@arg{#4}%
  \ifx\@fc@arg\empty
    % No separate bold font
    \setCJKmainfont{#3}[Path=#2,#1]%
    \setCJKsansfont{#3}[Path=#2,#1]%
    \setCJKmonofont{#3}[Path=#2,#1]%
  \else
    \setCJKmainfont{#3}[Path=#2,BoldFont=#4]%
    \setCJKsansfont{#3}[Path=#2,BoldFont=#4]%
    \setCJKmonofont{#3}[Path=#2,BoldFont=#4]%
  \fi
}

\newcommand{\@fc@warn}[1]{%
  \PackageWarningNoLine{font-config}{#1}%
}

% ── Priority 1: User Override ───────────────────────────────────────────────
\ifx\CJKFontPath\undefined
  % No user override — proceed to auto-detect
\else
  \@fc@warn{Using user-defined CJK font: \CJKFontPath\CJKFontName}
  \ifx\CJKFontBold\undefined
    \@fc@setfonts{\CJKFontPath}{\CJKFontName}{}%
  \else
    \@fc@setfonts{\CJKFontPath}{\CJKFontName}{\CJKFontBold}%
  \fi
  % Skip rest of file
  \expandafter\endinput
\fi

% ── Priority 1.5: Project-local .fonts/ — Source Han Serif SC ────────────────
% Ships as individual OTFs in .fonts/ (gitignored). Prefer Medium weight for
% body text on slides — heavier than the system .ttc Regular.
\IfFileExists{./.fonts/SourceHanSerifSC-Medium.otf}{%
  \@fc@warn{Detected project-local Source Han Serif SC (Medium+Bold)}
  \@fc@setfonts{./.fonts/}{SourceHanSerifSC-Medium.otf}{SourceHanSerifSC-Bold.otf}%
  \expandafter\endinput
}{}

% ── Priority 2: Windows — Microsoft YaHei ───────────────────────────────────
\IfFileExists{C:/WINDOWS/Fonts/msyh.ttc}{%
  \@fc@warn{Detected Windows — using Microsoft YaHei}
  \@fc@setfonts{C:/WINDOWS/Fonts/}{msyh.ttc}{msyhbd.ttc}%
}{%
% ── Priority 3: Linux — Noto Sans CJK (system package) ──────────────────────
  \IfFileExists{/usr/share/fonts/opentype/noto/NotoSansCJK-Regular.ttc}{%
    \@fc@warn{Detected Linux — using Noto Sans CJK (opentype)}
    \@fc@setfonts{/usr/share/fonts/opentype/noto/}{NotoSansCJK-Regular.ttc}{NotoSansCJK-Bold.ttc}%
  }{%
    \IfFileExists{/usr/share/fonts/truetype/noto/NotoSansSC-Regular.ttf}{%
      \@fc@warn{Detected Linux — using Noto Sans SC (truetype)}
      \@fc@setfonts{/usr/share/fonts/truetype/noto/}{NotoSansSC-Regular.ttf}{NotoSansSC-Bold.ttf}%
    }{%
% ── Priority 4: macOS — Noto Sans CJK or system fonts ───────────────────────
      \IfFileExists{/System/Library/Fonts/PingFang.ttc}{%
        \@fc@warn{Detected macOS — using PingFang}
        \@fc@setfonts{/System/Library/Fonts/}{PingFang.ttc}{}%
      }{%
        \IfFileExists{~/Library/Fonts/NotoSansSC-Regular.ttf}{%
          \@fc@warn{Detected macOS — using user-installed Noto Sans SC}
          \@fc@setfonts{~/Library/Fonts/}{NotoSansSC-Regular.ttf}{NotoSansSC-Bold.ttf}%
        }{%
% ── Priority 5: Generic fallback — search common paths ──────────────────────
          \IfFileExists{/usr/local/share/fonts/NotoSansSC-Regular.ttf}{%
            \@fc@warn{Using generic Noto Sans SC}
            \@fc@setfonts{/usr/local/share/fonts/}{NotoSansSC-Regular.ttf}{NotoSansSC-Bold.ttf}%
          }{%
% ── FAILURE: No CJK font found ──────────────────────────────────────────────
            \@fc@warn{No CJK font found on system. CJK text will not render correctly.}
            \@fc@warn{Install fonts: sudo apt-get install fonts-noto-cjk  (Linux)}
            \@fc@warn{Or set manually: \def\CJKFontPath{/path/} \def\CJKFontName{font.ttf}}
          }%
        }%
      }%
    }%
  }%
}
\makeatother


% ─── Packages ───────────────────────────────────────────────────────────────
\usepackage{booktabs,tabularx,multirow,array,microtype,colortbl}
\usepackage{graphicx}
\usepackage{xcolor}
\usepackage{adjustbox}
\usepackage{pifont}
\usepackage{tikz}
\usetikzlibrary{positioning,calc}

\graphicspath{{slides_assets/}}

% ═════════════════════════════════════════════════════════════════════════════
%  COLOR PALETTE — Navy + Gold
% ═════════════════════════════════════════════════════════════════════════════
\definecolor{Navy}{HTML}{1a3a6b}
\definecolor{NavyDark}{HTML}{0d1f3e}
\definecolor{NavyLight}{HTML}{c8d8ee}
\definecolor{NavyTint}{HTML}{e8eff8}
\definecolor{Brick}{HTML}{c9a23a}
\definecolor{BrickLight}{HTML}{fdf5dc}
\definecolor{Ink}{HTML}{1a1a2e}
\definecolor{InkSoft}{HTML}{5a5a7e}
\definecolor{Surface}{HTML}{f7f8fb}
\definecolor{PosGreen}{HTML}{2E7D32}
\definecolor{NegRed}{HTML}{C62828}

% ═════════════════════════════════════════════════════════════════════════════
%  BEAMER COLOR WIRING
% ═════════════════════════════════════════════════════════════════════════════
\setbeamercolor{background canvas}{bg=white}
\setbeamercolor{normal text}{fg=Ink,bg=white}
\setbeamercolor{frametitle}{bg=Navy,fg=white}
\setbeamercolor{frametitle right}{bg=NavyDark}
\setbeamercolor{title separator}{fg=Brick}
\setbeamercolor{progress bar}{fg=Brick,bg=NavyLight}
\setbeamercolor{alerted text}{fg=Brick}
\setbeamercolor{block title}{bg=NavyTint,fg=Navy}
\setbeamercolor{block body}{bg=Surface}
\setbeamercolor{block title alerted}{bg=BrickLight,fg=Brick}
\setbeamercolor{block body alerted}{bg=white}
\setbeamercolor{block title example}{bg=NavyLight,fg=NavyDark}
\setbeamercolor{block body example}{bg=white}
\setbeamercolor{palette primary}{fg=white,bg=Navy}
\setbeamercolor{itemize item}{fg=Navy}
\setbeamercolor{itemize subitem}{fg=Navy}
\setbeamercolor{enumerate item}{fg=Navy}
\setbeamercolor{enumerate subitem}{fg=Navy}

% ═════════════════════════════════════════════════════════════════════════════
%  BEAMER FONT SETTINGS
% ═════════════════════════════════════════════════════════════════════════════
\setbeamerfont{frametitle}{size=\Large,series=\bfseries}
\setbeamerfont{framesubtitle}{size=\normalsize}
\setbeamerfont{block title}{size=\normalsize,series=\bfseries}
\setbeamerfont{block body}{size=\small}

% ═════════════════════════════════════════════════════════════════════════════
%  NAVIGATION & FOOTLINE
% ═════════════════════════════════════════════════════════════════════════════
\setbeamertemplate{navigation symbols}{}
\setbeamertemplate{footline}{%
  \leavevmode%
  \hbox{%
    \begin{beamercolorbox}[wd=0.72\paperwidth,ht=2.5ex,dp=1.25ex,leftskip=8pt]{palette primary}%
      \usebeamerfont{footline}\tiny CvG-Diff MICCAI 2025%
    \end{beamercolorbox}%
    \begin{beamercolorbox}[wd=0.28\paperwidth,ht=2.5ex,dp=1.25ex,rightskip=8pt]{palette primary}%
      \hfill\usebeamerfont{footline}\tiny\insertframenumber\,/\,\inserttotalframenumber%
    \end{beamercolorbox}%
  }%
}

% ═════════════════════════════════════════════════════════════════════════════
%  CUSTOM COMMANDS
% ═════════════════════════════════════════════════════════════════════════════
\newcommand{\cmark}{\textcolor{PosGreen}{\ding{51}}}
\newcommand{\xmark}{\textcolor{NegRed}{\ding{55}}}
\newcommand{\deltapos}[1]{\textcolor{PosGreen}{\textbf{+#1}}}
\newcommand{\deltaneg}[1]{\textcolor{NegRed}{\textbf{#1}}}
\newcommand{\cnum}[1]{{\notocjk#1}}
\newcommand{\thetaEMA}{\theta^{\mathrm{EMA}}}

% ── Highlight box for key takeaways (lightweight, no heavy borders) ──────────
\newcommand{\keytakeaway}[1]{%
  \vspace{0.4em}%
  \begin{tikzpicture}
    \node[fill=BrickLight, rounded corners=3pt, inner sep=8pt, text width=\linewidth-16pt, align=justify] {%
      \small\textcolor{Brick}{\textbf{Key Takeaway.}} #1%
    };
  \end{tikzpicture}%
  \vspace{0.2em}%
}

% ── Inline highlight for important terms ─────────────────────────────────────
\newcommand{\hterm}[1]{\textcolor{Brick}{\textbf{#1}}}

% ── Auto image helper ────────────────────────────────────────────────────────
\newcommand{\autoimg}[2][]{%
  \centering\adjincludegraphics[max width=\linewidth,max height=0.72\textheight,#1]{#2}%
}

% ═════════════════════════════════════════════════════════════════════════════
%  METADATA
% ═════════════════════════════════════════════════════════════════════════════
\title{Cross-view Generalized Diffusion Model\\for Sparse-view CT Reconstruction}
\subtitle{CvG-Diff · MICCAI 2025}
\author{Jixiang Chen · Yiqun Lin · Yi Qin · Hualiang Wang · Xiaomeng Li}
\institute{HKUST xmed-lab \quad|\quad 复旦大学 FARM 组}
\date{2026年5月20日}

% ═════════════════════════════════════════════════════════════════════════════
\begin{document}
% ═════════════════════════════════════════════════════════════════════════════

% ── S1: Title ────────────────────────────────────────────────────────────────
{
\setbeamercolor{background canvas}{bg=NavyDark}
\begin{frame}[plain]
  \begin{columns}[T]
    \begin{column}{0.58\textwidth}
      \vspace{2em}
      {\color{Brick}\rule{0.7\linewidth}{1.5pt}}\\[0.6em]
      {\Large\bfseries\color{white}
        Cross-view Generalized Diffusion Model\\[0.2em]
        for Sparse-view CT Reconstruction}\\[0.6em]
      {\color{NavyLight}\large CvG-Diff · MICCAI 2025}\\[1em]
      {\small\color{NavyLight}
        Jixiang Chen · Yiqun Lin · Yi Qin\\[0.15em]
        Hualiang Wang · Xiaomeng Li}\\[0.3em]
      {\footnotesize\color{InkSoft} HKUST xmed-lab}\\[1.5em]
      {\color{Brick}\rule{0.7\linewidth}{0.5pt}}\\[0.6em]
      {\small\color{NavyLight} 2026年5月20日}
    \end{column}
    \begin{column}{0.40\textwidth}
      \vspace{0.5em}
      \includegraphics[width=\linewidth,height=0.70\textheight,keepaspectratio]{paper_fancy_title.jpg}
    \end{column}
  \end{columns}
\end{frame}
}

% ── S2: CT Imaging Background ────────────────────────────────────────────────
\begin{frame}{CT 成像原理}
  \vspace{0.3em}
  \begin{columns}[T]
    \begin{column}{0.55\textwidth}
      \textbf{三步成像流程}
      \begin{enumerate}
        \item \textbf{多角度投影}：X 射线从 $N_v$ 个角度穿过人体
        \item \textbf{Sinogram}：记录每个角度的衰减投影
        \item \textbf{滤波反投影（FBP）}：$\mathcal{A}^\dagger$ 将 sinogram 反投影重建 CT 图像
      \end{enumerate}
      \vspace{0.5em}
      \textbf{关键术语}
      \begin{itemize}
        \item CT（Computed Tomography）
        \item NFE（Number of Function Evaluations）
        \item Radon 变换 $\mathcal{A}$（正向投影）
        \item FBP 算子 $\mathcal{A}^\dagger$（反投影重建）
      \end{itemize}
    \end{column}
    \begin{column}{0.43\textwidth}
      \centering
      \adjincludegraphics[max width=0.48\textwidth,max height=0.35\textheight]{sparse-view-ct-1.jpg}%
      \hspace{4pt}%
      \adjincludegraphics[max width=0.48\textwidth,max height=0.35\textheight]{sparse-view-ct-2.jpg}
      \\[3pt]
      \footnotesize\color{InkSoft} CT sinogram 采集与 FBP 重建示意
    \end{column}
  \end{columns}
\end{frame}

% ── S3: Radon & FBP ──────────────────────────────────────────────────────────
\begin{frame}{Radon 变换与 FBP：CT 重建的正反算子}
  \vspace{0.3em}
  \begin{columns}[T]
    \begin{column}{0.48\textwidth}
      \textbf{Radon 变换 $\mathcal{A}$（正向）}
      \[
        \mathcal{A}f(\theta,s)=\int_{-\infty}^{+\infty}\!
        f(s\cos\theta - t\sin\theta,\;
          s\sin\theta + t\cos\theta)\,dt
      \]
      \vspace{0.3em}
      将图像 $f(x,y)$ 映射为 sinogram：$(\theta,s)$ 平面上每点是沿对应直线的线积分。
    \end{column}
    \begin{column}{0.48\textwidth}
      \textbf{滤波反投影 FBP（$\mathcal{A}^\dagger$）}
      \[
        f(x,y)=\int_0^{\pi}\!\bigl[
          \mathcal{A}f(\theta,s)*h(s)
        \bigr]_{s=x\cos\theta+y\sin\theta}\,d\theta
      \]
      \vspace{0.3em}
      $h(s)$：Ram-Lak 斜坡滤波器。视图数 $N_v$ 不足时，FBP 产生严重 \hterm{streak 伪影}。
    \end{column}
  \end{columns}

  \keytakeaway{%
    稀疏视图 CT 的核心挑战：如何从角度严重欠采样的 sinogram 重建诊断级图像？%
  }
\end{frame}

% ── S4: Sparse-view Problem ──────────────────────────────────────────────────
\begin{frame}{稀疏视图 CT：少剂量的代价}
  \vspace{0.3em}
  \begin{columns}[T]
    \begin{column}{0.60\textwidth}
      \centering
      \begin{tabular}{@{}c@{\hspace{6pt}}c@{\hspace{6pt}}c@{}}
        \adjincludegraphics[max width=0.30\textwidth,max height=0.38\textheight]{ours_recon_72v.png} &
        \adjincludegraphics[max width=0.30\textwidth,max height=0.38\textheight]{ours_recon_36v.png} &
        \adjincludegraphics[max width=0.30\textwidth,max height=0.38\textheight]{ours_recon_18v.png} \\[3pt]
        \small\textbf{72-view} & \small\textbf{36-view} & \small\textbf{18-view}\\
        \footnotesize 轻微伪影 & \footnotesize 明显 streak & \footnotesize 严重失真
      \end{tabular}
      \\[3pt]
      \footnotesize\color{InkSoft} FBP 直接重建结果（AAPM-LDCT 数据集）
    \end{column}
    \begin{column}{0.38\textwidth}
      \textbf{动机}
      \begin{itemize}
        \item View 数 $\downarrow$ $\Rightarrow$ X 射线剂量 $\downarrow$
        \item 目标：以\hterm{极少投影角度}重建\hterm{诊断级}图像
      \end{itemize}
      \vspace{0.5em}
      \textbf{挑战：FBP streak 伪影}
      \begin{itemize}
        \item View 越少 $\Rightarrow$ 欠采样越严重
        \item 18-view 条件下重建图像伪影严重，难以满足临床诊断要求
      \end{itemize}
    \end{column}
  \end{columns}
\end{frame}

% ── S5: Existing Methods ─────────────────────────────────────────────────────
\begin{frame}{现有方法：各有局限}
  \vspace{0.3em}
  \begin{columns}[T]
    \begin{column}{0.48\textwidth}
      \textbf{\cnum{①} 一步前馈方法（NFE = 1）}
      \\[2pt]
      \small\textbf{代表：}FreeSeed · GloReDi · DuDoTrans
      \begin{itemize}
        \item[\cmark] 推理效率高，NFE = 1，临床应用潜力
        \item[\xmark] 需\textbf{分别训练}，跨角度泛化差
        \item[\xmark] 极稀疏时\textbf{过平滑}，细节丢失
      \end{itemize}
    \end{column}
    \begin{column}{0.48\textwidth}
      \textbf{\cnum{②} 扩散模型方法}
      \\[2pt]
      \small\textbf{代表：}VSS · CoSIGN · DPS
      \begin{itemize}
        \item[\cmark] 生成质量高，可跨 view 数泛化
        \item[\xmark] NFE 高达\textbf{1000}，推理开销极大
        \item[\xmark] 极稀疏（18-view）时不稳定
      \end{itemize}
    \end{column}
  \end{columns}

  \keytakeaway{%
    目标：重建质量优 $+$ 低 NFE $+$ 跨视图泛化 $\;\Longrightarrow\;$ \textbf{CvG-Diff}%
  }
\end{frame}

% ── S6: Diffusion Primer ─────────────────────────────────────────────────────
\begin{frame}{预备知识：扩散模型}
  \vspace{0.3em}
  \begin{columns}[T]
    \begin{column}{0.52\textwidth}
      \textbf{标准 DDPM（Ho et al.~NeurIPS 2020）}
      \\[4pt]
      \textbf{前向加噪：}
      \[
        x_t=\sqrt{\bar\alpha_t}\,x_0+\sqrt{1-\bar\alpha_t}\,\varepsilon,
        \quad\varepsilon\sim\mathcal{N}(0,I)
      \]
      \vspace{0.3em}
      \small\textbf{反向去噪：}$\epsilon_\theta$ 迭代预测噪声
      \\
      \small\textbf{采样：}$x_T\!\sim\!\mathcal{N}(0,I)$ 迭代 $T$ 步
    \end{column}
    \begin{column}{0.45\textwidth}
      \textbf{Cold Diffusion（Bansal et al.~NeurIPS 2023）}
      \\[4pt]
      将加噪推广为\hterm{任意确定性退化算子} $D$：
      \begin{itemize}
        \item $D(x_0,0)=x_0$（无退化）
        \item $D(x_0,T)$：最严重退化
        \item 退化可以是模糊、掩码、或\textbf{稀疏视图子采样}
      \end{itemize}
    \end{column}
  \end{columns}

  \keytakeaway{%
    CvG-Diff：$D$ 定义为稀疏视图 CT 的角度子采样与 FBP%
  }
\end{frame}

% ── S7: Cross-view Observation ───────────────────────────────────────────────
\begin{frame}{关键观察：跨视图退化相关性}
  \vspace{0.3em}
  \begin{columns}[T]
    \begin{column}{0.52\textwidth}
      \centering
      \begin{tabular}{@{}c@{\hspace{4pt}}c@{\hspace{4pt}}c@{}}
        \adjincludegraphics[max width=0.30\textwidth,max height=0.35\textheight]{ours_recon_18v.png} &
        \adjincludegraphics[max width=0.30\textwidth,max height=0.35\textheight]{ours_recon_36v.png} &
        \adjincludegraphics[max width=0.30\textwidth,max height=0.35\textheight]{ours_recon_72v.png} \\[3pt]
        \scriptsize\textbf{18-view} & \scriptsize\textbf{36-view} & \scriptsize\textbf{72-view}\\
        \scriptsize 退化最重 & \scriptsize 中等 & \scriptsize 较轻
      \end{tabular}
      \\[3pt]
      \footnotesize\color{InkSoft} FBP 重建：基础解剖结构相同，退化程度各异
    \end{column}
    \begin{column}{0.46\textwidth}
      \textbf{核心洞察}
      \begin{itemize}
        \item 18/36/72-view 均来自同一 sinogram 子采样
        \item View 越少 $\Rightarrow$ 退化越重
        \item 基础解剖结构\textbf{完全一致}
      \end{itemize}
      \vspace{0.5em}
      同一解剖结构在不同稀疏率下 = 同一图像的\hterm{不同退化级别}
    \end{column}
  \end{columns}

  \keytakeaway{%
    退化是\textbf{确定性的} $\Rightarrow$ Cold Diffusion 框架直接适用 $\Rightarrow$ \textbf{单一模型}即可泛化至所有投影角度数%
  }
\end{frame}

% ── S8: Cold Diffusion Framework ─────────────────────────────────────────────
\begin{frame}{Cold Diffusion 框架（Bansal et al.~NeurIPS 2023）}
  \vspace{0.3em}
  \begin{columns}[T]
    \begin{column}{0.50\textwidth}
      \textbf{训练目标}
      \[
        \mathcal{L}_{\mathrm{restore}}
        =\bigl\|R_\theta(D(x_0,t),t)-x_0\bigr\|_2^2
      \]
      \vspace{0.3em}
      \footnotesize 随机取 $t\in\{1,\ldots,T\}$，退化到 level $t$，网络还原
      \vspace{0.8em}

      \textbf{推理：迭代更新}
      \begin{align*}
        \hat{x}_0^{\,t} &= R_\theta(x_t,t)\\
        x_{t-1} &= x_t - D(\hat{x}_0^{\,t},t)+D(\hat{x}_0^{\,t},t-1)
      \end{align*}
    \end{column}
    \begin{column}{0.47\textwidth}
      \textbf{框架性质}
      \begin{itemize}
        \item $D$ \textbf{不限于高斯噪声}，任意确定性退化均可
        \item 边界：$D(x_0,0)=x_0$，$D(x_0,T)$=最严重退化
        \item 推理链：$x_T\!\to\!x_{T-1}\!\to\!\cdots\!\to\!x_0$
      \end{itemize}
      \vspace{0.5em}
      \textbf{CvG-Diff 适配}
      \begin{itemize}
        \item $D$ 定义为 \textbf{稀疏视图子采样 + FBP}
        \item $x_T$ 含真实测量，残差为数据一致性约束
      \end{itemize}
    \end{column}
  \end{columns}
\end{frame}

% ── S9: CvG-Diff Framework ───────────────────────────────────────────────────
\begin{frame}{CvG-Diff 框架}
  \vspace{0.5em}
  \centering
  \adjincludegraphics[max width=0.85\textwidth,max height=0.55\textheight]{paper_fig1_framework.png}
  \\[0.8em]
  \begin{columns}[T,onlytextwidth]
    \begin{column}{0.58\textwidth}
      \textbf{退化算子 $D$}
      \\[2pt]
      \small $x_T=\mathcal{A}^\dagger P(T)\,\mathcal{A}\,x_0$
      \quad\footnotesize（$\mathcal{A}$：Radon；$P(T)$：子采样；$\mathcal{A}^\dagger$：FBP）
    \end{column}
    \begin{column}{0.38\textwidth}
      \textbf{实验关键级别}
      \\[2pt]
      \centering\small Level\,5/7/8 $\;\Leftrightarrow\;$ 72/36/18 view
    \end{column}
  \end{columns}
\end{frame}

% ── S10: Error Propagation ───────────────────────────────────────────────────
\begin{frame}{直接应用 Cold Diffusion 的失效模式：误差传播}
  \vspace{0.5em}
  \centering
  \adjincludegraphics[max width=0.90\textwidth,max height=0.48\textheight]{paper_fig2_error_propagation.png}
  \\[0.8em]
  \begin{columns}[T,onlytextwidth]
    \begin{column}{0.55\textwidth}
      \textbf{误差累积机制}
      \\[2pt]
      $\hat{x}_0^{\,t}$ 含还原误差 $\Rightarrow$ 误差注入 $x_{t-1}$ $\Rightarrow$ 后续输入质量持续下降
      \\[2pt]
      $\Rightarrow$ \hterm{误差逐步累积，无法逆转}
    \end{column}
    \begin{column}{0.42\textwidth}
      \textbf{训练-推理分布偏移}
      \\[2pt]
      $R_\theta$ 在训练阶段从未接触含误差的中间状态输入
      \\[2pt]
      $\Rightarrow$ \hterm{训练-推理分布偏移}
    \end{column}
  \end{columns}
\end{frame}

% ── S11: EPCT ────────────────────────────────────────────────────────────────
\begin{frame}{对策一（训练侧）：EPCT — 主动模拟误差传播}
  \vspace{0.3em}
  \begin{columns}[T]
    \begin{column}{0.52\textwidth}
      \textbf{EPCT 核心思路}
      \\[3pt]
      训练阶段\hterm{主动合成}推理过程中将遭遇的含误差中间状态——使网络具备对带传播误差输入的\textbf{鲁棒还原能力}
      \vspace{0.5em}

      \textbf{Step 1}：维护网络的 EMA 版本 $R_{\thetaEMA}$
      \[
        \thetaEMA\;\leftarrow\;
        \gamma\,\thetaEMA+(1-\gamma)\,\theta
      \]
      \vspace{0.2em}
      \footnotesize $\gamma=0.995$，每 $p=10$ 次 iter 更新一次
      \vspace{0.5em}

      \textbf{Step 2}：EMA 在最严重退化级别 $T$ 预测
      \[
        \hat{x}_0^{\,T}=R_{\thetaEMA}(x_T,\,T)
      \]
    \end{column}
    \begin{column}{0.45\textwidth}
      \textbf{合成带传播误差的 $x_t$}
      \\[3pt]
      随机取 $t\in[1,T)$，令
      \[
        x_t=x_T-D(\hat{x}_0^{\,T},T)+D(\hat{x}_0^{\,T},t)
      \]
      \vspace{0.2em}
      \footnotesize $x_t$ 模拟推理第 $t$ 步实际遇到的\textbf{带误差输入}
      \vspace{0.8em}

      \textbf{实现细节}
      \begin{itemize}
        \item EMA 提供\textbf{稳定、滞后}的版本，避免 bootstrapping 不稳定
        \item EPCT 在训练迭代 $>$ 2000 步后启用
      \end{itemize}
    \end{column}
  \end{columns}
\end{frame}

% ── S12: EPCT Dual Loss ──────────────────────────────────────────────────────
\begin{frame}{对策一（训练侧）：EPCT — 双 Loss 联合训练}
  \vspace{0.3em}
  \begin{columns}[T]
    \begin{column}{0.50\textwidth}
      \textbf{原始 Restore Loss}
      \[
        \mathcal{L}_{\mathrm{restore}}
        =\bigl\|R_\theta(D(x_0,t),t)-x_0\bigr\|_2^2
      \]
      \vspace{0.2em}
      \footnotesize 从\textbf{理想退化图像}还原，维持基础能力
      \vspace{0.8em}

      \textbf{新增 Compose Loss}
      \[
        \mathcal{L}_{\mathrm{compose}}
        =\bigl\|x_0-R_\theta(x_t,t)\bigr\|_2^2
      \]
      \vspace{0.2em}
      \footnotesize 从\textbf{带传播误差}的 $x_t$ 反演，提升鲁棒性
    \end{column}
    \begin{column}{0.47\textwidth}
      \textbf{为什么有效？}
      \begin{itemize}
        \item $\mathcal{L}_{\mathrm{restore}}$：从\textbf{理想}退化输入重建（覆盖第一步）
        \item $\mathcal{L}_{\mathrm{compose}}$：从\textbf{含误差}输入重建（覆盖后续）
        \item 二者\textbf{互补}，填补训练-推理分布偏移
      \end{itemize}
      \vspace{0.5em}
      \textbf{消融验证}
      \\[3pt]
      \centering\small
      仅 EPCT $\Rightarrow$ 18-view \deltapos{4.18 dB}\\
      \textbf{（主要贡献模块）}
    \end{column}
  \end{columns}

  \keytakeaway{%
    $\mathcal{L}=\mathcal{L}_{\mathrm{restore}}+\mathcal{L}_{\mathrm{compose}}$%
  }
\end{frame}

% ── S13: Sequential Sampling Limitations ─────────────────────────────────────
\begin{frame}{残留问题：顺序采样的两个局限}
  \vspace{0.3em}
  \begin{columns}[T]
    \begin{column}{0.48\textwidth}
      \textbf{\cnum{①} 语义错误在稀疏级被锁死}
      \\[3pt]
      稀疏 level（大 $t$）产生的\hterm{语义错误}（如解剖结构位置偏差、结构缺失）
      \\[3pt]
      固定顺序 $T\!\to\!T\!-\!1\!\to\!\cdots\!\to\!1$：语义错误在后续 dense level 步骤\textbf{无法修复}
      \vspace{0.8em}

      \textbf{\cnum{②} NFE 分配浪费}
      \\[3pt]
      固定顺序调度：每个 level 各消耗 1 NFE
      \\[2pt]
      $\Rightarrow$ 某 level 已收敛仍被迫继续
      \\[2pt]
      $\Rightarrow$ NFE 未必花在最需要的地方
    \end{column}
    \begin{column}{0.49\textwidth}
      \textbf{核心观察}
      \\[3pt]
      稀疏级的语义错误 $\neq$ dense 级能修复
      \\[3pt]
      必须先在\hterm{语义层面}（稀疏 level）充分校正，再进入高频细化
      \vspace{1em}

      \begin{alertblock}{SPDPS 动机}
        $\Rightarrow$ 需要\textbf{自适应、双阶段}的采样策略
        \\[4pt]
        这正是 \textbf{SPDPS} 的设计动机。
      \end{alertblock}
    \end{column}
  \end{columns}
\end{frame}

% ── S14: SPDPS ───────────────────────────────────────────────────────────────
\begin{frame}{对策二（推理侧）：SPDPS — 语义优先双阶段采样}
  \vspace{0.3em}
  \centering
  \adjincludegraphics[max width=0.90\textwidth,max height=0.48\textheight]{paper_fig4_spdps_nocap.png}
  \\[0.5em]
  \footnotesize\color{InkSoft} 上行：SPDPS 重采样策略；下行：顺序采样基线
  \vspace{0.5em}

  \begin{columns}[T,onlytextwidth]
    \begin{column}{0.48\textwidth}
      \textbf{Phase~1：语义校正（$n=6$）}
      \\[3pt]
      \centering
      若 $\mathrm{SSIM}(\hat{x}_0^{\,t},\hat{x}_0^{\,t+1})>\tau=0.97$
      \\[3pt]
      \textbf{重置}至 $T_{\max}$，继续语义校正
    \end{column}
    \begin{column}{0.48\textwidth}
      \textbf{Phase~2：高频细化（$m=4$）}
      \\[3pt]
      \centering
      Phase~1 完成后
      \\[3pt]
      dense level 顺序 $m$ 步采样
    \end{column}
  \end{columns}
\end{frame}

% ── S15: Quantitative Results ────────────────────────────────────────────────
\begin{frame}{定量评估：各视图数均优于最优基线}
  \vspace{0.5em}
  {\centering\footnotesize
  \setlength{\tabcolsep}{5pt}%
  \begin{tabular}{@{}l r r r r@{}}
    \toprule
    & & \multicolumn{3}{c}{\textbf{PSNR\,$\uparrow$\;/\;SSIM\,$\uparrow$\quad（AAPM-LDCT）}} \\[-2pt]
    \cmidrule(l){3-5}
    \textbf{方法} & \textbf{NFE} & \textbf{18-view} & \textbf{36-view} & \textbf{72-view} \\
    \midrule
    DuDoTrans  & 1    & 34.02\;/\;.901 & 38.24\;/\;.941 & 42.76\;/\;.976 \\
    FreeSeed   & 1    & 34.31\;/\;.904 & 37.92\;/\;.939 & 42.93\;/\;.975 \\
    GloReDi    & 1    & 34.75\;/\;.913 & 37.54\;/\;.935 & 41.80\;/\;.963 \\
    \midrule
    VSS        & 1000 & 32.34\;/\;.879 & 37.52\;/\;.940 & 41.92\;/\;.971 \\
    CoSIGN     & 10   & 31.84\;/\;.863 & 34.96\;/\;.897 & 37.87\;/\;.932 \\
    \midrule
    \rowcolor{NavyLight}
    \textbf{CvG-Diff} & \textbf{10} & \textbf{38.34\;/\;.952} & \textbf{41.78\;/\;.971} & \textbf{45.94\;/\;.986} \\
    \bottomrule
  \end{tabular}\par}

  \keytakeaway{%
    18-view 领先最佳一步法 \deltapos{+3.6 dB}，领先 VSS \deltapos{+6.0 dB}；NFE 仅需 \textbf{10}（vs VSS 的 1000）%
  }
\end{frame}

% ── S16: Qualitative Results ─────────────────────────────────────────────────
\begin{frame}{视觉质量：18-view 重建对比}
  \vspace{0.5em}
  \centering
  \adjincludegraphics[max width=0.92\textwidth,max height=0.62\textheight]{paper_fig3_qualitative.png}

  \keytakeaway{%
    CvG-Diff（最右列）：streak 伪影得到有效消除，高频细节保真度最高；PSNR 领先最优基线 \deltapos{6 dB}%
  }
\end{frame}

% ── S17: Ablation ────────────────────────────────────────────────────────────
\begin{frame}{消融实验：EPCT 为主要贡献模块}
  \vspace{0.5em}
  \centering
  \adjincludegraphics[max width=0.80\textwidth,max height=0.50\textheight]{ablation_bar.png}

  \keytakeaway{%
    18-view PSNR 消融：仅 SPDPS \deltapos{+1.00 dB} ／ 仅 EPCT \deltapos{+4.18 dB} ／ +SPDPS \deltapos{+0.49 dB}%
  }
\end{frame}

% ── S18: Ablation Interpretation ─────────────────────────────────────────────
\begin{frame}{消融实验：解读}
  \vspace{0.3em}
  \begin{columns}[T]
    \begin{column}{0.40\textwidth}
      \centering\small
      \begin{tabular}{lc}
        \toprule
        \textbf{变体} & \textbf{增益}\\
        \midrule
        仅 SPDPS   & \deltapos{1.00 dB}\\
        仅 EPCT    & \deltapos{4.18 dB}\\
        +SPDPS     & \deltapos{0.49 dB}\\
        \bottomrule
      \end{tabular}
    \end{column}
    \begin{column}{0.56\textwidth}
      \textbf{物理解释}
      \begin{itemize}
        \item EPCT：消除\textbf{训练-推理分布偏移}
        \item SPDPS：解决\textbf{语义错误锁死}问题
        \item 双模块协同改善同一误差传播瓶颈
      \end{itemize}
    \end{column}
  \end{columns}

  \keytakeaway{%
    \textbf{关键洞察}：EPCT 核心贡献 \deltapos{+4.18 dB}（训练侧）$+$ SPDPS 推理优化 \deltapos{+0.49 dB}（推理侧）
    两者互补增益叠加，\textbf{双管齐下}消除误差传播%
  }
\end{frame}

% ── S19: Reproduction ────────────────────────────────────────────────────────
\begin{frame}{复现：数值对比}
  \vspace{0.3em}
  \centering\footnotesize
  \setlength\tabcolsep{4pt}
  \begin{tabular}{lcccccc}
    \toprule
    & \multicolumn{2}{c}{\textbf{18-view}} &
      \multicolumn{2}{c}{\textbf{36-view}} &
      \multicolumn{2}{c}{\textbf{72-view}}\\
    \cmidrule(lr){2-3}\cmidrule(lr){4-5}\cmidrule(lr){6-7}
    \textbf{方法} & PSNR & SSIM & PSNR & SSIM & PSNR & SSIM\\
    \midrule
    Paper Table 1      & 38.34  & 0.9518 & 41.78  & 0.9705 & 45.94  & 0.9863\\
    HF 官方 ckpt       & 38.335 & 0.9518 & 41.782 & 0.9705 & 45.670 & 0.9854\\
    \rowcolor{NavyLight}
    \textbf{本组（从头训）} &
      \textbf{38.403} & \textbf{0.9539} &
      \textbf{41.856} & \textbf{0.9722} &
      \textbf{45.666} & \textbf{0.9863}\\
    \midrule
    $\Delta$ 本组 $-$ Paper &
      \deltapos{0.063} & \deltapos{0.0021} &
      \deltapos{0.076} & \deltapos{0.0017} &
      \deltaneg{$-$0.274} & $\pm$0.000\\
    \bottomrule
  \end{tabular}

  \keytakeaway{%
    72-view $-$0.27 dB 与 HF 官方 ckpt 一致 $\Rightarrow$ torch-radon 在 CUDA 12.4 重编的\textbf{数值微差}，非训练问题%
  }
\end{frame}

% ── S20: Qualitative Reproduction ────────────────────────────────────────────
\begin{frame}{复现：定性可视化（18-view，L506 患者）}
  \vspace{0.3em}
  \centering
  \adjincludegraphics[max width=0.92\textwidth,max height=0.55\textheight]{ours_recon_18v.png}
  \\[0.5em]
  \footnotesize\color{InkSoft} 列：GT / FBP / HF ckpt / 本组
  \vspace{0.3em}

  \begin{columns}[T,onlytextwidth]
    \begin{column}{0.58\textwidth}
      \textbf{视觉观察}
      \\[2pt]
      \small
      FBP 列严重 streak 伪影；HF 官方 ckpt 与\textbf{本组结果}视觉差异\textbf{可忽略不计}；Streak 伪影经 SPDPS 有效消除，软组织边界与参考图像高度吻合
    \end{column}
    \begin{column}{0.38\textwidth}
      \begin{alertblock}{定性确认}
        \small 从头训练结果与官方 ckpt \textbf{等价}
      \end{alertblock}
    \end{column}
  \end{columns}
\end{frame}

% ── S21: Closing ─────────────────────────────────────────────────────────────
{
\setbeamercolor{background canvas}{bg=NavyDark}
\begin{frame}[plain]
  \vfill
  \centering
  {\color{white}\fontsize{36}{44}\selectfont 谢谢聆听}\\[0.5em]
  {\color{Brick}\large Questions Welcome}\\[2em]
  {\color{NavyLight}\normalsize
    CvG-Diff · MICCAI 2025 · arXiv 2508.10313}\\[0.3em]
  {\color{InkSoft}\small
    Jixiang Chen · Yiqun Lin · Yi Qin · Hualiang Wang · Xiaomeng Li\\[0.2em]
    HKUST xmed-lab}
  \vfill
\end{frame}
}

% ── Backup Slides ────────────────────────────────────────────────────────────
\appendix

\begin{frame}{Backup：SPDPS 超参敏感度（Table 3）}
  \vspace{0.5em}
  \centering
  \adjincludegraphics[max width=0.85\textwidth,max height=0.60\textheight]{paper_table3_tau_m.png}
  \\[0.5em]
  \small\centering
  $\tau\in\{0.97,0.98,0.99\}$ 差 $<$0.4 dB；$m\in\{2,3,4\}$ 效果几乎相同。
\end{frame}

\begin{frame}{Backup：36-view 定性结果}
  \vspace{0.5em}
  \centering
  \adjincludegraphics[max width=0.92\textwidth,max height=0.75\textheight]{ours_recon_36v.png}
\end{frame}

\begin{frame}{Backup：72-view 定性结果}
  \vspace{0.5em}
  \centering
  \adjincludegraphics[max width=0.92\textwidth,max height=0.75\textheight]{ours_recon_72v.png}
\end{frame}

\end{document}
